\section{A failed search exposes a shortage}
For $S\subseteq L$, let $N(S)$ be the tasks acceptable to at least one student in $S$. Suppose a complete alternating search from an unmatched student $x$ reaches no unused task. Let $A$ be the reached students, including $x$, and let $B$ be the reached tasks.

Before the general counting argument, consider three students $a,b,c$ who can all use only tasks $1,2$. Suppose $a$ is matched to $1$ and $b$ to $2$, with root $x=c$. The search reaches both tasks and then their owners $a,b$. It finds no unused task. Thus $A=\{c,a,b\}$ and $B=\{1,2\}$.

The existing matching pairs task $1$ with $a$ and task $2$ with $b$. Every reached student \emph{except the unmatched root} appears in exactly one such pair. There is no pair for $c$, which is precisely the shortage. We now show that this one-extra-student pattern holds whenever a complete search fails, not just in this example.

Every task in $B$ is matched; otherwise the search would have succeeded. Its matched student is reached by the search rule and lies in $A\setminus\{x\}$. Conversely, every reached student other than $x$ was reached from a matched task. Because matching partners are unique, the matching pairs give a bijection between $B$ and $A\setminus\{x\}$. Therefore
\[
|A|=|B|+1.
\]

A bijection here is not an abstract new object that we must discover. It is the existing partner rule: send each task in $B$ to its owner. Two tasks cannot have the same owner because $M$ is a matching. Every student in $A\setminus\{x\}$ is an owner of a reached task because that is the only rule that can reach a new student. Thus the partner rule is injective and surjective onto that set. Equal finite sizes follow, and adding back the single root gives the displayed equation.

The separate identity $N(A)=B$ is needed because the shortage condition counts \emph{all acceptable tasks} of $A$, not merely whichever tasks the search happened to record. A proof that reached tasks are too few would be inconclusive if some acceptable tasks had been missed. The next two-containment argument rules that out.

We also have $N(A)=B$. Every reached task was entered from a reached student, so $B\subseteq N(A)$. For the reverse containment, take a task neighboring a reached student. If their connecting edge is unmatched, the exploration follows it, so the task lies in $B$. If their edge is matched, the student is not the unmatched root $x$, and its matched task is the task from which it was reached; that task already lies in $B$. Thus all neighbors lie in $B$.

It helps to run this argument on the three-student example above, where $a$ holds task $1$, $b$ holds task $2$, the root is $c$, and every student accepts only tasks $1$ and $2$. For the root $c$, both edges $c$--$1$ and $c$--$2$ are unmatched, so the search follows them and both tasks lie in $B$. For $a$, the edge $a$--$1$ is matched, and task $1$ is exactly the task from which $a$ was reached; the edge $a$--$2$ is unmatched, so the search follows it, and task $2$ is in $B$ anyway. The same reasoning applies to $b$. Every acceptable task of every reached student is in $B=\{1,2\}$, as the general argument promised. If the search had stopped early, say after reaching only task $1$ and its owner $a$, then task $2$ would be an acceptable task of a reached student lying outside the recorded set, and the comparison of sizes below would prove nothing. This is why the search must be complete.

Combining the two identities gives
\[
|N(A)|=|B|=|A|-1<|A|.
\]
There are fewer available tasks than students in $A$. No assignment can give all those students distinct acceptable tasks. The search has produced an independently checkable certificate of impossibility.

\begin{principle}{Habit 7: make failure informative}
A failed attempt is not automatically a counterexample. Ask whether the failure can be turned into a concrete object that violates a necessary condition. In this problem, the object is a set of students competing for too few tasks. In another problem it may be two distinct inputs with the same output, refuting a claim of injectivity, or a particular value violating a required inequality.
\end{principle}
